
\documentclass{article}
\usepackage{colortbl}
\usepackage{makecell}
\usepackage{multirow}
\usepackage{supertabular}

\begin{document}

\newcounter{utterance}

\centering \large Interaction Transcript for game `cladder', experiment `full\_v1.5\_default', episode 2646 with llama{-}3.1{-}8b{-}instruct.
\vspace{24pt}

{ \footnotesize  \setcounter{utterance}{1}
\setlength{\tabcolsep}{0pt}
\begin{supertabular}{c@{$\;$}|p{.15\linewidth}@{}p{.15\linewidth}p{.15\linewidth}p{.15\linewidth}p{.15\linewidth}p{.15\linewidth}}
    \# & \multicolumn{2}{c}{Player} && \multicolumn{2}{c}{Game Master} \\
    \hline

    \theutterance \stepcounter{utterance}  
    & & & \multicolumn{4}{p{0.6\linewidth}}{
        \cellcolor[rgb]{0.9,0.9,0.9}{
            \makecell[{{p{\linewidth}}}]{
                \texttt{\tiny{[P1$\langle$GM]}}
                \texttt{You are an expert in causal inference. The following question is not a typical commonsense query, but rather a meticulously designed question created by a professor specializing in causal inference, intended to assess the students' mastery of the course content.} \\
\texttt{Imagine a self{-}contained, hypothetical world with only the following conditions, and without any unmentioned factors or causal relationships: Gender has a direct effect on smoking and lung cancer. Smoking has a direct effect on tar deposit. Tar deposit has a direct effect on lung cancer. Gender is unobserved. For nonsmokers and with no tar deposit, the probability of lung cancer is 79\%. For nonsmokers and with high tar deposit, the probability of lung cancer is 49\%. For smokers and with no tar deposit, the probability of lung cancer is 71\%. For smokers and with high tar deposit, the probability of lung cancer is 40\%. For nonsmokers, the probability of high tar deposit is 79\%. For smokers, the probability of high tar deposit is 40\%. For smokers, would it be more likely to see lung cancer if the person had been a nonsmoker?} \\
\texttt{Start your answer with 'Yes' or 'No', followed by additional reasoning or evidence to support your explanation.} \\
            }
        }
    }
     \\ \\

    \theutterance \stepcounter{utterance}  
    & \multicolumn{4}{p{0.6\linewidth}}{
        \cellcolor[rgb]{0.561,0.737,0.561}{
            \makecell[{{p{\linewidth}}}]{
                \texttt{\tiny{[P1$\rangle$GM]}}
                \texttt{No.} \\
\\ 
\texttt{To address this question, we need to apply the principles of causal inference, specifically the concept of counterfactuals. We are asked to compare the probability of lung cancer in smokers who were previously nonsmokers (the treatment group) with the probability of lung cancer in smokers who were always smokers (the control group).} \\
\\ 
\texttt{Given the information provided, we can use the potential outcomes framework to estimate the counterfactual probability of lung cancer in smokers who were previously nonsmokers. Let's denote the potential outcomes as follows:} \\
\\ 
\texttt{{-} Y(0) = Lung cancer in a nonsmoker} \\
\texttt{{-} Y(1) = Lung cancer in a smoker} \\
\\ 
\texttt{We are given the following probabilities:} \\
\\ 
\texttt{{-} P(Y(0) = 1 | no tar deposit) = 0.79} \\
\texttt{{-} P(Y(0) = 1 | high tar deposit) = 0.49} \\
\texttt{{-} P(Y(1) = 1 | no tar deposit) = 0.71} \\
\texttt{{-} P(Y(1) = 1 | high tar deposit) = 0.40} \\
\texttt{{-} P(high tar deposit | nonsmoker) = 0.79} \\
\texttt{{-} P(high tar deposit | smoker) = 0.40} \\
\\ 
\texttt{We can use the law of total probability to estimate the counterfactual probability of lung cancer in smokers who were previously nonsmokers:} \\
\\ 
\texttt{P(Y(1) = 1 | previously nonsmoker) = P(high tar deposit | previously nonsmoker) \textbackslash{}* P(Y(1) = 1 | high tar deposit) + P(no tar deposit | previously nonsmoker) \textbackslash{}* P(Y(1) = 1 | no tar deposit)} \\
\\ 
\texttt{However, we don't have direct information about the probability of tar deposit in smokers who were previously nonsmokers. To address this, we can use the fact that the probability of tar deposit is 40\% in smokers who were always smokers. We can assume that the probability of tar deposit in smokers who were previously nonsmokers is similar to that in smokers who were always smokers, given that smoking is the common cause of tar deposit in both groups.} \\
\\ 
\texttt{Therefore, we can estimate the counterfactual probability of lung cancer in smokers who were previously nonsmokers as follows:} \\
\\ 
\texttt{P(Y(1) = 1 | previously nonsmoker) ≈ P(high tar deposit | smoker) \textbackslash{}* P(Y(1) = 1 | high tar deposit) + P(no tar deposit | smoker) \textbackslash{}* P(Y(1) = 1 | no tar deposit)} \\
\texttt{= 0.40 \textbackslash{}* 0.40 + 0.60 \textbackslash{}* 0.71} \\
\texttt{= 0.16 + 0.426} \\
\texttt{= 0.586} \\
\\ 
\texttt{Now, we can compare this estimate with the probability of lung cancer in smokers who were always smokers, which is 0.40. We can see that the estimated probability of lung cancer in smokers who were previously nons} \\
            }
        }
    }
    & & \\ \\

    \theutterance \stepcounter{utterance}  
    & & & \multicolumn{2}{p{0.3\linewidth}}{
        \cellcolor[rgb]{0.9,0.9,0.9}{
            \makecell[{{p{\linewidth}}}]{
                \texttt{\tiny{[GM$|$GM]}}
                \texttt{No} \\
            }
        }
    }
    & & \\ \\

    \theutterance \stepcounter{utterance}  
    & & & \multicolumn{2}{p{0.3\linewidth}}{
        \cellcolor[rgb]{0.9,0.9,0.9}{
            \makecell[{{p{\linewidth}}}]{
                \texttt{\tiny{[GM$|$GM]}}
                \texttt{no} \\
            }
        }
    }
    & & \\ \\

    \theutterance \stepcounter{utterance}  
    & & & \multicolumn{2}{p{0.3\linewidth}}{
        \cellcolor[rgb]{0.9,0.9,0.9}{
            \makecell[{{p{\linewidth}}}]{
                \texttt{\tiny{[GM$|$GM]}}
                \texttt{game\_result = WIN} \\
            }
        }
    }
    & & \\ \\

\end{supertabular}
}

\end{document}
